\section*{अध्याय \ref{ch:b2:wittig} --- विटिग और C=C बनाने वाली अन्य अभिक्रियाएँ}

\begin{solution}{exo:b2:wittig:1}
\ce{Ph3P + CH3CH2Br -> [Ph3PCH2CH3]+ + Br-} (द्वि-अणुक प्रतिस्थापन, किसी विलायक में गर्म करके), फिर
अक्रिय गैस के अंतर्गत शुष्क ईथर या THF में ब्यूटिललिथियम: \ce{[Ph3PCH2CH3]+ + BuLi -> Ph3P=CHCH3 + BuH + Li+}।
\end{solution}

\begin{solution}{exo:b2:wittig:2}
इलाइड \ce{Ph3P=CH2} और बेन्ज़ैल्डिहाइड फ़ेनिलएथीन (स्टाइरीन), \ce{C6H5CH=CH2}, और
ट्राइफ़ेनिलफ़ॉस्फ़ीन ऑक्साइड देते हैं।
\end{solution}

\begin{solution}{exo:b2:wittig:3}
\ce{Ph3P=CHCOOEt} और \ce{Ph3P=CHCOCH3}: स्थायीकृत। \ce{Ph3P=CHCH3}: अस्थायीकृत, वही जिसे
ब्यूटिललिथियम (या सोडियम हाइड्राइड, सोडियम ऐमाइड) चाहिए। \ce{Ph3P=CHC6H5}: ऐरिल-स्थायीकृत, व्यवहार में किसी
ऐल्कॉक्साइड या हाइड्रॉक्साइड से बनाया जाता है, और ज्यामितियों का मिश्रण देता है।
\end{solution}

\begin{solution}{exo:b2:wittig:4}
एथिल (\textit{E})-पेंट-2-ईनोएट, \ce{CH3CH2CH=CHCOOEt}, साथ में सोडियम डाइएथिल फ़ॉस्फ़ेट।
\end{solution}

\begin{solution}{exo:b2:wittig:5}
विटिग: मेथिलट्राइफ़ेनिलफ़ॉस्फ़ोनियम आयोडाइड और ब्यूटिललिथियम से बना इलाइड \ce{Ph3P=CH2} साइक्लोहेक्सेनोन के साथ केवल
मेथिलीनसाइक्लोहेक्सेन देता है, द्वि-आबंध वलय के बाहर। 1-मेथिलसाइक्लोहेक्सेन-1-ऑल का निर्जलीकरण
एक तृतीयक कार्बधनायन से होकर जाता है और वह हाइड्रोजन खोता है जो अधिक प्रतिस्थापित ऐल्कीन देता है (ज़ैत्सेव):
मुख्यतः 1-मेथिलसाइक्लोहेक्सीन, ग़लत समावयवी।
\end{solution}

\begin{solution}{exo:b2:wittig:6}
दोनों कटाव एक ही हैं: बेन्ज़ैल्डिहाइड और बेन्ज़िलिडीन इलाइड \ce{Ph3P=CHC6H5}। वह इलाइड केवल
ऐरिल-स्थायीकृत है और (\textit{E})/(\textit{Z}) मिश्रण देता है। शुद्ध (\textit{E}) के लिए बेन्ज़ैल्डिहाइड के साथ
डाइएथिल बेन्ज़िलफ़ॉस्फ़ोनेट की HWE अभिक्रिया प्रयुक्त कीजिए।
\end{solution}

\begin{solution}{exo:b2:wittig:7}
\ce{Ph3P=CHCH3} के साथ: (\textit{Z})-पेंट-2-ईन, \ce{CH3CH2CH=CHCH3}, मुख्यतः (\textit{Z})।
\ce{Ph3P=CHCOOEt} के साथ: एथिल (\textit{E})-पेंट-2-ईनोएट, मुख्यतः (\textit{E})।
\end{solution}

\begin{solution}{exo:b2:wittig:8}\emergencystretch=8em
\ce{Ph3P=CH2 + C6H10O -> C7H12 + Ph3PO}। द्रव्यमान: मेथिलीनसाइक्लोहेक्सेन \qty{96.173}{g/mol},
ट्राइफ़ेनिलफ़ॉस्फ़ीन ऑक्साइड \qty{278.291}{g/mol}; परमाणु मितव्ययिता $96.173/(96.173 + 278.291) =
96.173/374.464 \approx 25.7~\%$। अभिकर्मकों के द्रव्यमान का तीन-चौथाई फ़ॉस्फ़ीन ऑक्साइड के रूप में निकल जाता है।
% ledger: aw:C, aw:H, aw:O, aw:P
\end{solution}

\begin{solution}{exo:b2:wittig:9}
नाभिकरागी फ़ॉस्फ़ोरस ब्रोमाइड को विस्थापित करता है:
\begin{equation*}
\ce{P(OEt)3 + BrCH2COOEt -> [(EtO)3PCH2COOEt]+ + Br-}.
\end{equation*}
फिर ब्रोमाइड
फ़ॉस्फ़ोनियम आयन के किसी एथिल समूह के एक कार्बन पर आक्रमण करता है (द्वि-अणुक प्रतिस्थापन), जिससे
ब्रोमोएथेन मुक्त होता है और \ce{P=O} आबंध बनता है:
\begin{equation*}
\ce{[(EtO)3PCH2COOEt]+ + Br- -> (EtO)2P(O)CH2COOEt + EtBr}.
\end{equation*}
\end{solution}

\begin{solution}{exo:b2:wittig:10}
\ce{OHC-CH=CH-CHO + 2Ph3P=CHC6H5 -> C6H5CH=CH-CH=CH-CH=CHC6H5 + 2Ph3PO}: बेन्ज़िलिडीन इलाइड के दो तुल्यांक
(बेन्ज़िलट्राइफ़ेनिलफ़ॉस्फ़ोनियम क्लोराइड और किसी क्षारक से), हर ऐल्डिहाइड समूह के लिए एक।
नए द्वि-आबंध ज्यामितियों के मिश्रण के रूप में बनते हैं; पूर्णतः-(\textit{E}) समावयवी, जो सबसे अधिक
स्थायी है, क्रिस्टलन द्वारा पृथक किया जाता है, या मिश्रण का समावयवीकरण किया जाता है।
\end{solution}

\begin{solution}{exo:b2:wittig:11}
\omterm{def:b2:enolates-aldol:aldol}{ऐल्डोल}: आधिक्य प्रोपेनोन और सोडियम हाइड्रॉक्साइड के साथ बेन्ज़ैल्डिहाइड (\cref{ch:b2:enolates-aldol});
सस्ते अभिकर्मक, उपोत्पाद जल, (\textit{E}) उत्पाद, परंतु प्रोपेनोन आधिक्य में न हो तो संघनन दो बार हो सकता है
(डाइबेन्ज़ैल्ऐसीटोन)। विटिग: \omterm{def:b2:wittig:ylide}{स्थायीकृत इलाइड}
\ce{Ph3P=CHCOCH3} के साथ बेन्ज़ैल्डिहाइड; (\textit{E}), एक ही उत्पाद, कोई दोहरा संघनन नहीं, परंतु इलाइड महँगा है और
ट्राइफ़ेनिलफ़ॉस्फ़ीन ऑक्साइड हटाना पड़ता है। यहाँ \omterm{def:b2:enolates-aldol:aldol}{ऐल्डोल} बेहतर मार्ग है।
\end{solution}

\begin{solution}{exo:b2:wittig:12}
विटिग: लवण-मुक्त परिस्थितियों में \ce{Ph3P=CHCH2CH2CH3} (1-ब्रोमोब्यूटेन से, फिर ब्यूटिललिथियम) के साथ ब्यूटेनल:
मुख्यतः (\textit{Z}); उपोत्पाद ट्राइफ़ेनिलफ़ॉस्फ़ीन ऑक्साइड, और कुछ (\textit{E}) समावयवी। ऐल्काइन:
ऑक्ट-4-आइन (एथाइन से, 1-ब्रोमोप्रोपेन और सोडियम ऐमाइड के साथ दो ऐल्किलनों द्वारा) का
विषाक्त पैलेडियम उत्प्रेरक पर हाइड्रोजनन (\cref{prop:b2:alkene-redox:syn-hydrogenation}): केवल (\textit{Z});
उपोत्पाद ऐल्किलनों से सोडियम ब्रोमाइड और अमोनिया हैं।
\end{solution}

\begin{solution}{pb:b2:wittig:1}\emergencystretch=3em
\textbf{1.} \ce{C23H46}, अर्थात् $n = 23$ के साथ \ce{C_nH_{2n}}: एक द्वि-आबंध (या एक वलय), यहाँ एक
\ce{C=C}।
\textbf{2.} \ce{C9=C10} द्वि-आबंध।
\textbf{3.} 1-हैलोटेट्राडेकेन से बने इलाइड के साथ नोनेनल \ce{CH3(CH2)7CHO}; या 1-हैलोनोनेन से बने इलाइड के साथ टेट्राडेकेनल
\ce{CH3(CH2)12CHO}।
\textbf{4.} 1-ब्रोमोटेट्राडेकेन, \ce{CH3(CH2)13Br}।
\textbf{5.} भारी-भरकम ट्राइफ़ेनिलफ़ॉस्फ़ीन द्वारा द्वि-अणुक प्रतिस्थापन किसी अबाधित प्राथमिक कार्बन पर तेज़ है
और विलोपन से प्रतिस्पर्धा नहीं होती।
\textbf{6.} \ce{Ph3P + C14H29Br -> [Ph3PC14H29]+ + Br-}, टेट्राडेसिलट्राइफ़ेनिलफ़ॉस्फ़ोनियम ब्रोमाइड।
\textbf{7.} ब्यूटिललिथियम (या सोडियम हाइड्राइड, सोडियम ऐमाइड); अस्थायीकृत इलाइड के लिए हाइड्रॉक्साइड बहुत दुर्बल
क्षारक है, और उसके साथ आया जल इलाइड को प्रोटॉनित कर देता।
\textbf{8.} नहीं, उस पर केवल एक ऐल्किल समूह है: मुख्यतः (\textit{Z}) (\cref{prop:b2:wittig:stereo})।
\textbf{9.} एक चतुःसदस्यी वलय \ce{P-C-C-O}: फ़ॉस्फ़ोरस उस कार्बन से आबंधित है जिस पर
\ce{C13H27} शृंखला है, वह कार्बन उस कार्बन से जिस पर \ce{C8H17} शृंखला है, और वह ऑक्सीजन से, जो
फ़ॉस्फ़ोरस से आबंधित है।
\textbf{10.} ब्यूटिललिथियम और इलाइड जल और ऑक्सीजन से अभिक्रिया करते हैं।
\textbf{11.} द्वि-आबंध नोनेनल के भूतपूर्व कार्बोनिल कार्बन (उत्पाद का C9) को भूतपूर्व
इलाइड कार्बन (C10) से जोड़ता है, और कोई अन्य आबंध नहीं खिसकता (\cref{prop:b2:wittig:regiospecific})।
\textbf{12.} C9 पर एक कार्बधनायन, जो C8 या C10 से हाइड्रोजन खो सकता है और खिसक सकता है:
ट्राइकोस-8-ईन और ट्राइकोस-9-ईन का मिश्रण, हर एक (\textit{E}) और (\textit{Z}) के रूप में, अधिकांशतः (\textit{E})।
\textbf{13.} \ce{Ph3P=CHC13H27 + C8H17CHO -> C8H17CH=CHC13H27 + Ph3PO}।
\textbf{14.} लक्ष्य \ce{C23H46}: \qty{322.621}{g/mol}; नोनेनल \ce{C9H18O}: \qty{142.242}{g/mol}; लवण
\ce{C32H44BrP}: \qty{539.582}{g/mol}; \ce{Ph3PO}, \ce{C18H15OP}: \qty{278.291}{g/mol}।
\textbf{15.} इलाइड \ce{C32H43P} का द्रव्यमान $539.582 - 80.912 = \qty{458.670}{g/mol}$ है; नोनेनल के साथ
अभिकर्मकों के \qty{600.912}{g/mol} से उत्पाद के \qty{322.621}{g/mol}: $322.621/600.912 \approx
53.7~\%$।
\textbf{16.} इसकी ध्रुवता से: अध्रुवी हाइड्रोकार्बन अध्रुवी निक्षालक के साथ सिलिका के एक छोटे स्तंभ से
निकल जाता है जबकि ऑक्साइड रुका रहता है; या ऑक्साइड किसी ठंडे अध्रुवी विलायक से क्रिस्टलित होकर अलग हो जाता है।
\textbf{17.} बनने वाला अत्यंत प्रबल \ce{P=O} आबंध टूटने वाले \ce{P-C} आबंध से भारी पड़ता है
(\cref{prop:b2:wittig:driving})।
\textbf{18.} इसके लिए ऐसा \omterm{def:b2:wittig:hwe}{फ़ॉस्फ़ोनेट} चाहिए जिसके अभिक्रियाशील कार्बन पर इलेक्ट्रॉन-आकर्षी समूह हो, जो
साधारण ऐल्किल शृंखला में नहीं है, और यह (\textit{E})-ऐल्कीन देती है।
\textbf{19.} ट्राइकोस-9-आइन, \ce{CH3(CH2)7C#C(CH2)12CH3}, किसी विषाक्त पैलेडियम
उत्प्रेरक (लेड लवण और क्विनोलीन के साथ कैल्सियम कार्बोनेट पर पैलेडियम) पर हाइड्रोजन के साथ।
\textbf{20.} डेक-1-आइन, \ce{CH3(CH2)7C#CH}, सोडियम ऐमाइड द्वारा विप्रोटॉनित, फिर 1-ब्रोमोट्राइडेकेन
\ce{CH3(CH2)12Br}: $8 + 2 + 13 = 23$ कार्बन, त्रि-आबंध C9 और C10 के बीच।
\textbf{21.} आगे का हाइड्रोजनन ट्राइकोसेन देता, जिसमें कोई द्वि-आबंध नहीं; विषाक्त उत्प्रेरक
ऐल्काइन को ऐल्कीन से कहीं अधिक प्रबलता से बाँधता है, और ऐल्कीन दूसरे योग से पहले सतह छोड़ देती है।
\textbf{22.} दोनों उपलब्ध टुकड़ों से दो चरण लेते हैं। ऐल्काइन मार्ग (\textit{Z}) समावयवी
स्वच्छ रूप से देता है और केवल लवण छोड़ता है; विटिग मार्ग कुछ (\textit{E}) समावयवी और
ट्राइफ़ेनिलफ़ॉस्फ़ीन ऑक्साइड का बड़ा द्रव्यमान देता है।
\textbf{23.} 100~\%: \ce{C23H44 + H2 -> C23H46}, हर परमाणु उत्पाद में पहुँचता है।
\textbf{24.} फ़ेरोमोन के प्रति ग्राम $278.291/322.621 = \boldsymbol{\approx \qty{0.86}{g}}$
ट्राइफ़ेनिलफ़ॉस्फ़ीन ऑक्साइड।
\end{solution}
